\clearpage
\section{Nhiệm vụ 2 -- Phương pháp mới và biến đổi đồ thị}
Nhiệm vụ 2 được tổ chức thành hai phần: Task A đáp ứng yêu cầu dataset thứ hai
và phương pháp từ năm 2023 trở lại đây; Task B học thao tác biến đổi đồ thị.
GQA thuộc Task A, không thuộc chương VG150.

\begin{figure}[H]\centering
\begin{tikzpicture}[font=\small,>=Latex,
b/.style={draw=inkblue,rounded corners,fill=mutedblue,align=center,minimum height=10mm,text width=63mm}]
\node[b,text width=55mm] (root) {Nhiệm vụ 2};
\node[b,anchor=north] (a) at (-3.8,-1.2) {Task A: dataset thứ hai + phương pháp mới\\Custom GQA: MLP/GAT/PE-Net/AlignG};
\node[b,anchor=north] (b) at (3.8,-1.2) {Task B: graph transformation\\MagicBrush $\to$ EMNLP 2020};
\draw[->] (root.south)--(a.north);
\draw[->] (root.south)--(b.north);
\end{tikzpicture}
\caption{Hai phần của Nhiệm vụ 2. EMNLP là dataset thay thế ở Task B.}
\end{figure}

\subsection{Task A -- Custom GQA và phương pháp mới}
GQA có scene graph chuẩn hóa, phù hợp tái sử dụng pipeline PredCls với
đặc trưng CLIP đóng băng~\cite{hudson2019gqa}. Bản custom dùng 69.942 ảnh
train, 5.000 ảnh validation và 310 predicate foreground; đây không phải
benchmark GQA200. Các mô hình so sánh dùng cùng split và cùng evaluator nội bộ.

\begin{table}[H]\centering\small
\caption{Vai trò của từng cấu hình Task A.}\label{tab:gqamethod}
\begin{tabularx}{\linewidth}{@{}>{\raggedright\arraybackslash}p{26mm}X@{}}\toprule
Cấu hình & Thành phần được kiểm tra\\\midrule
MLP / GAT & Baseline phi đồ thị và message passing.\\
PE-CE & Nhánh prototype theo PE-Net, huấn luyện với cross-entropy.\\
PE-PL & Thêm prototype learning (PL).\\
PE-PLPR & Thêm prototype regularization (PR) vào PE-PL.\\
PE-PLPR + AlignG & Thêm cơ chế context/feedback vào baseline PE-PLPR trong matched ablation.\\
\bottomrule\end{tabularx}\end{table}

PE-Net~\cite{zheng2023penet} được adapted cho feature đã cache:
đặc trưng cặp thay cho union-box feature của detector gốc, có khởi tạo GloVe.
Nhánh AlignG~\cite{jung2026aligng} bổ sung prototype context và GRU feedback.
Do feature, label space và pipeline khác bản gốc, đây là \textbf{triển khai
adapted}, không phải reproduction hay so sánh trực tiếp với số trong paper.

\reporttodo{2A-1}{Bổ sung sơ đồ tensor của bản adapted, phương trình loss PL/PR,
trọng số từng loss và bảng ánh xạ module với paper từ source/config đã khóa.
Không điền công thức theo paper nếu code thực tế đã thay đổi.}
\reporttodo{2A-2}{Hoàn thiện data card custom GQA: cách lọc vocabulary, manifest chia ảnh,
và trạng thái held-out test. Các bảng hiện tại là validation; chưa dùng làm kết quả test.}
\srcnote{Nguồn: \texttt{docs/LEVEL2\_METHOD\_DECISION.md};
\texttt{reports/level2/gqa-full-v3-results-v1/README.md}.}

\clearpage
\subsection{Task A -- Bảng kết quả và matched ablation}
\begin{table}[H]\centering\small
\caption{Custom GQA validation; GC, mean $\pm$ sample std (\%), seed 42/123/3407.}\label{tab:gqa}
\begin{tabular}{@{}lrrrr@{}}\toprule
Model & R@50 & mR@50 & R@100 & mR@100\\\midrule
MLP & $76,66\pm0,59$ & $3,33\pm0,34$ & $87,24\pm0,40$ & $5,68\pm0,75$\\
GAT & $77,34\pm0,13$ & $3,10\pm0,10$ & $87,83\pm0,12$ & $4,84\pm0,22$\\
PE-CE & $78,07\pm0,27$ & $3,00\pm0,28$ & $88,39\pm0,25$ & $4,88\pm0,58$\\
PE-PL & $\mathbf{78,14\pm0,07}$ & $3,42\pm0,14$ & $\mathbf{88,51\pm0,05}$ & $5,61\pm0,48$\\
PE-PLPR & $77,97\pm0,07$ & $\mathbf{4,16\pm0,17}$ & $88,28\pm0,08$ & $\mathbf{7,13\pm0,31}$\\
\bottomrule\end{tabular}\end{table}

R@50 quanh 77--78\% nhưng mR@50 chỉ 3--4\%: recall tổng thể che khuất chênh lệch
giữa predicate. PE-PLPR tăng 1,16 pp mR@50 so PE-CE, trong khi R@50
giảm 0,10 pp. GAT tăng R nhưng giảm mR so MLP, khác với xu hướng trên VG150.

\begin{table}[H]\centering\small
\caption{Matched ablation AlignG trên cùng custom GQA validation.}\label{tab:aligng}
\begin{tabular}{@{}lrrr@{}}\toprule
Model & GC R@50 & GC mR@50 & NG mR@50\\\midrule
PE-PLPR & $77,97\pm0,07$ & $4,16\pm0,17$ & $7,72\pm0,26$\\
PE-PLPR + AlignG & $77,95\pm0,16$ & $4,59\pm0,39$ & $9,12\pm1,11$\\
Delta từ mean hiển thị (pp) & $-0,02$ & +0,43 & +1,40\\
\bottomrule\end{tabular}\end{table}

\begin{figure}[H]\centering
\begin{tikzpicture}
\begin{axis}[width=.85\linewidth,height=4.7cm,ybar,bar width=25pt,
ymin=-.25,ymax=1.65,xtick={1,2,3},
xticklabels={GC R@50,GC mR@50,NG mR@50},
ylabel={$\Delta$ AlignG (pp)},ymajorgrids,
tick label style={font=\small},axis x line*=bottom,axis y line*=left]
\addplot[fill=inkblue] coordinates {(1,-.02)(2,.43)(3,1.40)};
\draw[dashed] (axis cs:.5,0)--(axis cs:3.5,0);
\end{axis}\end{tikzpicture}
\caption{AlignG: thay đổi của mean hiển thị so với PE-PLPR. Không dùng ba metric
làm ba lần đo độc lập; độ biến động của từng model nằm trong bảng~\ref{tab:aligng}.}
\end{figure}

Matched ablation dùng cùng seed, batch 512, patience 15 và chọn checkpoint
theo validation GC mR@50. \textbf{PE-PLPR và AlignG cải thiện mean Recall,
không cải thiện đồng thời Recall tổng thể.}
\reporttodo{2A-3}{Thêm case đúng/sai có ảnh và graph, đối chiếu cùng cặp object
giữa PE-PLPR và AlignG. Chưa quy mức tăng mR cho một loại quan hệ cụ thể
khi chưa có phân tích per-predicate.}
\srcnote{Nguồn: \texttt{gqa-full-v3-results-v1/README.md};
\texttt{aligng-gqa-fast-v1-results-v1/summary.json} và \texttt{comparison.json}.
Delta tính từ số đầy đủ có thể lệch 0,01 pp so với hiệu hai mean đã làm tròn.}
